%=====================================================================
\chapter{What E-Commerce Is, and What It Changed}\label{ch:what}
%=====================================================================
\locator{1}{Part I --- E-Commerce and Its Models}

\begin{outcomes}
By the end of this chapter you will be able to:
\begin{tight}
  \item \textbf{Define} electronic commerce, and \textbf{distinguish} it from
        the wider category of electronic business.
  \item \textbf{Name} the four functions every shop must perform, and
        \textbf{locate} each of them in a storefront you have not seen before.
  \item \textbf{Explain} why selling online is a software problem rather than a
        marketing one, giving two consequences that follow.
  \item \textbf{Distinguish} the parts of a transaction that electronic commerce
        actually changed from the parts it only moved.
  \item \textbf{Classify} a given online business by the four functions, and
        \textbf{identify} which of them it has chosen not to perform itself.
\end{tight}
\end{outcomes}

\begin{prereq}
Nothing in this book. You need to have used an online shop as a customer, which
is the only prerequisite this chapter has. Appendix~A sets up the software you
will need from \chref{stack} onward; start it this week, because it sometimes
needs a second attempt.
\end{prereq}

\begin{keyterms}
electronic commerce $\bullet$ electronic business $\bullet$ storefront
$\bullet$ catalogue $\bullet$ basket $\bullet$ checkout $\bullet$ fulfilment
$\bullet$ transaction $\bullet$ disintermediation $\bullet$ long tail
$\bullet$ friction $\bullet$ conversion $\bullet$ abandonment
$\bullet$ trust
\end{keyterms}

\section{A Shop With No Closing Time}

Consider a bookshop on a street in Bahawalpur. It opens at nine and closes at
nine. It holds perhaps four thousand titles, because that is what the shelves
take. Its customers are the people who can reach the street. If you want a book
it does not stock, the shopkeeper can order it, and you will come back next
week.

Now consider the same bookshop online. It is open at three in the morning. It
can list four hundred thousand titles, because a listing is a database row and
rows are cheap. Its customers are anyone with a connection. If you want a book
it does not stock, the website can know that instantly, and can offer to tell
you when it arrives.

Every one of those differences is real, and none of them is the interesting one.

The interesting difference is this: in the street shop, a human being stands
between the customer and every mistake. If the price label is wrong, the
shopkeeper notices. If you try to buy two of something when one is left, the
shopkeeper tells you. If you hand over a torn note, the shopkeeper decides. The
shop runs on a continuous supply of human judgement, and that judgement is
invisible precisely because it never fails conspicuously.

Online, there is no shopkeeper. Every one of those judgements has to have been
written down, in advance, as a rule --- and anything that was not written down
will happen.

\begin{definitionbox}[Electronic Commerce]
\term{Electronic commerce} is the buying and selling of goods and services over
a computer network, together with the transfer of money and data that completes
those transactions.

\smallskip
It is a subset of \term{electronic business}, which is the use of networks for
\emph{any} business process --- a supplier portal, an internal stock system, a
staff payroll. The distinguishing feature of e-commerce is the
\term{transaction}: value changes hands, and both parties are entitled to
expect that it did so exactly once.

\smallskip
That last clause is not pedantry. It is the reason \chref{checkout} is the
longest chapter in this book.
\end{definitionbox}

\section{The Four Functions}

Strip away the branding and every shop on the internet does the same four
things, in the same order. This is worth fixing now, because it is the skeleton
of the whole course: each function is a part of this book, and each is a cluster
of screens in the software you will install in \chref{stack}.

\dsfig{%
\begin{tikzpicture}[
  fn/.style={rounded corners=3pt, draw=#1, line width=1.1pt, fill=#1!8,
             text=#1!60!black, font=\bfseries\small, align=center,
             minimum width=2.9cm, minimum height=1.25cm},
  arr/.style={-{Stealth[length=2.4mm]}, primary!55, line width=1pt},
  lbl/.style={font=\scriptsize\itshape, text=neutral, align=center}
]
  \node[fn=secondary] (c1) at (0,0)    {Catalogue\\{\scriptsize\mdseries what is
       for sale}};
  \node[fn=greenhl]   (c2) at (3.9,0)  {Basket\\{\scriptsize\mdseries what you
       chose}};
  \node[fn=purplehl]  (c3) at (7.8,0)  {Payment\\{\scriptsize\mdseries money
       changes hands}};
  \node[fn=tealhl]    (c4) at (11.7,0) {Fulfilment\\{\scriptsize\mdseries the
       thing arrives}};
  \draw[arr] (c1) -- (c2);
  \draw[arr] (c2) -- (c3);
  \draw[arr] (c3) -- (c4);

  \node[lbl] at (0,-1.15)    {Part II};
  \node[lbl] at (3.9,-1.15)  {Part III};
  \node[lbl] at (7.8,-1.15)  {Part IV};
  \node[lbl] at (11.7,-1.15) {Part V};

  \node[font=\scriptsize\itshape, text=accent, align=center] at (5.85,1.30)
       {reversible without consequence $\;\longrightarrow\;$ not reversible};
  \draw[dashed, accent!60, line width=0.8pt] (5.85,0.85) -- (5.85,-1.5);
\end{tikzpicture}}%
{The four functions of any shop, and the parts of this book that build them.
The dashed line is the one that matters: everything to its left can be undone
by closing the browser tab, and nothing to its right can.}{fourfunctions}

\begin{conceptbox}[The dashed line is the whole subject]
A customer browsing a catalogue and filling a basket has committed to nothing.
They can close the tab and the shop is exactly as it was. Mistakes on that side
of the line cost a sale at worst.

\smallskip
The moment payment is attempted, that stops being true. Money may have moved
while the network dropped. An order may exist that the customer does not know
about, or the customer may believe in an order that does not exist. Stock may
have been reserved for someone who never paid.

\smallskip
Most of what distinguishes a real shop from a student exercise is the handling
of that line, and most student projects put all their effort on the left of it
because that is the half you can see.
\end{conceptbox}

\section{What Actually Changed}

It is easy to overstate what e-commerce changed. A great deal of it is the same
commerce, conducted through a different window. Three things, though, genuinely
changed, and they explain most of what the rest of this book worries about.

\subsection{The cost of listing fell to nearly nothing}

Shelf space is expensive and finite; a database row is neither. A street shop
must stock what sells often, because an item that sells twice a year occupies
the same shelf as one that sells twice a week. An online shop has no such
pressure, and can profitably carry items that sell twice a year --- the
\term{long tail}.

That is a business observation with an engineering consequence that this course
will measure: a catalogue of four hundred thousand items does not behave like a
catalogue of four thousand. \chref{catalogue} puts fifty thousand products into
your shop and times the result, and the answer is not comfortable.

\subsection{Some middlemen became unnecessary, and others appeared}

\term{Disintermediation} is the removal of intermediaries: a manufacturer can
sell directly to you, and the distributor and the retailer are not needed.

The prediction made in the 1990s was that this would flatten commerce. What
happened instead was that old intermediaries were replaced by new and larger
ones --- marketplaces, payment processors, delivery networks, search engines.
A shop today typically pays a marketplace for access to customers, a gateway
for the ability to take money, and a courier to deliver, and each of those is a
dependency with a price and a failure mode. \chref{payments} is about one of
them in detail.

\subsection{Trust had to be manufactured}

In a street shop, trust is cheap. The shop is physically there; you can see the
goods; if it cheats you, you know where it is.

Online, none of that holds, and every substitute has to be built deliberately:
a visible address and a returns policy, a padlock in the address bar, reviews,
an order confirmation that arrives within seconds, a refund that works. These
are not decoration --- they are the mechanism that replaces a shopfront, and
each of them is a feature somebody has to implement and maintain.

\begin{examplebox}[Three signals, and what each costs to fake]
A customer deciding whether to buy from an unfamiliar shop is reading signals.
It is worth knowing which ones are hard to fake and which are not.

\begin{tight}
  \item \textbf{A professional-looking design.} Costs almost nothing to fake;
        a template is free. Carries little information.
  \item \textbf{A valid certificate and HTTPS.} Costs nothing now that
        certificates are free --- which is why its absence is damning and its
        presence proves very little. \chref{security} returns to exactly what
        it does and does not prove.
  \item \textbf{A confirmation email, with an order number, that arrives in
        four seconds.} Hard to fake, because it requires the shop to actually
        have the machinery. This is one reason \chref{checkout} treats the
        confirmation as part of the transaction rather than as an
        afterthought.
\end{tight}
\end{examplebox}

\section{Why This Is a Software Problem}

The common way to teach this subject is as business studies with screenshots.
That gets the emphasis wrong, and the published outline this course follows ---
see the Course Specification --- gets it right: eighteen of its twenty-two
clauses are about building, configuring and running the software.

Here is the argument in one paragraph. A shop is a set of rules that must hold
without supervision: this price applies to this customer in this quantity in
this region on this date; this voucher may be used once; this stock is reserved
for eleven minutes; this payment either completed or it did not. Every one of
those is a statement about state, every one of them has an edge, and nobody is
watching. That is not a marketing problem. It is the same problem as any other
concurrent, stateful, partially-trusted system, and it is solved with the same
tools.

\begin{alertbox}[The two failures that are specific to this subject]
Almost everything that goes badly wrong in an online shop is one of two things,
and both are invisible from the front page.

\smallskip
\textbf{The state got out of step.} Money taken, no order recorded. Order
recorded, stock never reduced. Refund issued twice. These are not rare, and they
do not announce themselves --- the shop looks fine, and the accounts do not
balance at the end of the month. \chref{checkout} and \chref{orders} are about
not causing them.

\smallskip
\textbf{Something was trusted that should not have been.} A price taken from a
form field rather than from the database. A quantity that was allowed to be
negative. An administrative page that was never actually checked for a login.
\chref{security} is a whole chapter about this, and its laboratory finds several
of them in the shop you will have built by then.
\end{alertbox}

\section{Reading a Shop You Have Not Seen Before}

The skill this chapter is really teaching is a diagnostic one, and it is worth
practising before you build anything. Given any online business, you should be
able to say within a few minutes which of the four functions it performs itself,
which it has handed to somebody else, and what that choice costs it.

\begin{worked}[Classifying three businesses by the four functions]
\textbf{A clothing brand selling from its own site.}
Catalogue: its own. Basket: its own. Payment: a gateway --- it does not hold a
banking licence, so this is handed over. Fulfilment: usually a courier.
\emph{Performs two of four; the two that define it.}

\smallskip
\textbf{The same brand selling only through a large marketplace.}
Catalogue: the marketplace's, in the marketplace's format. Basket: the
marketplace's. Payment: the marketplace's. Fulfilment: possibly the
marketplace's warehouse. \emph{Performs none of the four.} It has bought
access to customers and sold its relationship with them; it does not know who
they are, cannot email them, and can be delisted. The commission is the
visible price and is not the largest one.

\smallskip
\textbf{A food delivery platform.}
Catalogue: aggregated from restaurants. Basket: its own. Payment: its own.
Fulfilment: its own riders --- which is the expensive one, and the one that is
hard to copy. \emph{Performs three and a half of four}, and the half it does
not perform (the cooking) is the one with the thinnest margin.

\smallskip
\textbf{The general rule.} The function a business insists on performing itself
is usually the one it believes is its advantage. Ask which function a business
refuses to outsource, and you have found what it thinks it is.
\end{worked}

\begin{pitfall}
\begin{tight}
  \item \textbf{Confusing e-commerce with having a website.} A site that shows a
        catalogue and a phone number is a brochure. The transaction is the
        thing; without it, none of the hard problems in this book arise.
  \item \textbf{Assuming the hard part is the shopping cart.} It is not, and it
        has not been for twenty years --- you will configure one in an
        afternoon in \chref{basket}. The hard parts are everything after the
        dashed line in \dsref{fourfunctions}.
  \item \textbf{Treating trust signals as design.} A returns policy is not a
        page; it is a promise that somebody has to keep, and a refund path
        somebody has to build (\chref{orders}).
  \item \textbf{Believing disintermediation removed the middlemen.} It replaced
        them. Count your dependencies before claiming you have none.
  \item \textbf{Saying ``we will add payments later''.} Everything to the right
        of the dashed line constrains everything to its left. Deciding how you
        take money late is how projects discover their data model is wrong.
\end{tight}
\end{pitfall}

\begin{lab}[Reading four real shops]
No software is needed for this one. \chref{stack} installs the stack; this week
you are learning to look.

\smallskip
Choose \textbf{four} online businesses that you can actually reach, of
deliberately different shapes --- for example a single-brand store, a large
general marketplace, a services business that takes bookings rather than
shipping goods, and a local business selling through social media only.

\smallskip
For each one, as a customer and without buying anything, record:

\begin{tightnum}
  \item \textbf{The four functions.} Who performs each --- the business, or
        somebody else? How can you tell? For payment, the giveaway is usually
        whether the address bar changes during checkout.
  \item \textbf{Where the catalogue ends.} How many categories deep does it go?
        Can a product appear in two categories? Does it have options --- size,
        colour --- and are they separate products or one product with choices?
        You are looking at a data model through a window; \chref{catalogue} will
        show you the model itself.
  \item \textbf{The furthest point you can reach without an account.} Can you
        fill a basket? Reach the checkout? See the shipping cost before
        committing? Write down the exact step at which it stopped.
  \item \textbf{Three trust signals}, and for each, whether it is cheap or
        expensive to fake.
  \item \textbf{One thing that is broken or confusing.} There is always one.
        Describe it precisely enough that somebody could fix it.
\end{tightnum}

\smallskip
\textbf{Write up one page} comparing the four, ending with a single sentence
for each: what does this business think it is? Keep it --- \chref{planning}
asks you to do the same exercise for the shop you are going to build, and it is
much easier the second time.

\smallskip
\textbf{What to notice.} The four shops will differ far more in what they
\emph{refuse} to do than in what they do. Everyone has a catalogue. Not everyone
will tell you the delivery charge before you have handed over your address, and
the ones that will not have usually decided that on purpose.
\end{lab}

\begin{checkpoint}
\begin{tightnum}
  \item A university lets students view results and download transcripts
        online, but fee payment happens at a bank counter. Is this electronic
        commerce, electronic business, both or neither? Justify your answer from
        the definitions.
  \item Give one judgement a human shopkeeper makes without noticing, and
        describe precisely what has to exist in software to replace it.
  \item A marketplace charges a 15\% commission. A seller calculates that
        running its own site would cost less than 15\% of revenue and concludes
        it should leave. Name two costs the seller has not counted, and one
        thing it would gain that is not money.
\end{tightnum}
\end{checkpoint}

\begin{chaptersummary}
\begin{tight}
  \item Electronic commerce is the buying and selling of goods and services over
        a network, \emph{including} the transfer of money. The transaction is
        what distinguishes it from electronic business generally.
  \item Every shop performs four functions in order: catalogue, basket, payment,
        fulfilment. They are Parts II to V of this book.
  \item The line between basket and payment is the one that matters. Everything
        before it is reversible by closing a tab; nothing after it is.
  \item Three things genuinely changed: the cost of listing collapsed, old
        intermediaries were replaced by new and larger ones, and trust stopped
        being free and had to be manufactured.
  \item The absent shopkeeper is the engineering problem. Every judgement a
        human would have made has to have been written down in advance, and
        whatever was not written down will eventually happen.
  \item Almost all serious failures are one of two kinds: state that got out of
        step, or something trusted that should not have been.
  \item To read an unfamiliar business, ask which of the four functions it
        performs itself. The one it refuses to outsource is what it believes it
        is.
\end{tight}
\end{chaptersummary}

\begin{reviewq}
  \item \textbf{(MCQ)} Which of these is electronic business but \emph{not}
        electronic commerce? (a)~A customer buying a phone from a web shop;
        (b)~a supplier checking a manufacturer's stock levels through a portal;
        (c)~a customer paying a utility bill through an app; (d)~a marketplace
        taking commission on a sale.
  \item \textbf{(MCQ)} In \dsref{fourfunctions}, the dashed line separates:
        (a)~front end from back end; (b)~free actions from paid actions;
        (c)~what a customer can undo by closing the tab from what they cannot;
        (d)~catalogue data from order data.
  \item \textbf{(Short)} Explain the long tail, and give one engineering
        consequence of it that a shop owner would not anticipate.
  \item \textbf{(Short)} ``Disintermediation removed the middleman.'' Argue
        against this with a concrete example.
  \item \textbf{(Short)} A friend says the hard part of building an online shop
        is the shopping cart. Give the strongest version of why they are wrong,
        and name the two failure modes you would be more worried about.
  \item \textbf{(Applied)} A bookshop with four thousand titles in a shop wants
        to list eighty thousand online, most of which it does not hold in stock.
        Identify three decisions this forces that would not arise at four
        thousand titles, and say for each which of the four functions it
        affects.
  \item \textbf{(Applied)} Take any one of the four businesses from this
        chapter's laboratory. Describe, in enough detail that it could be
        checked, an experiment a customer could run from outside that would
        reveal whether the shop reserves stock during checkout. State what
        result would mean ``yes'', what would mean ``no'', and what would be
        inconclusive.
\end{reviewq}
